\chapter{Background and Related Work}
\label{ch:background}

\section{Economic games as behavioral probes}
\label{sec:bg-games}

A handful of small economic games have done outsized work in the study of cooperation,
because each isolates one tension and turns it into a number. They strip away the
particulars of any real situation and leave a controlled choice between acting for
oneself and acting for the group, played for real stakes and recorded move by move. That
is what is needed to ask whether one kind of agent exploits another: a setting where
exploitation, if it happens, shows up in the payoffs. Two such games anchor this thesis,
and a third property, iteration, makes them sharp enough to read.

\subsection{The Public Goods Game}
\label{subsec:bg-pgg}

In the Public Goods Game, each of $n$ players receives an endowment and privately chooses
how much of it to put into a common pool. The pool is multiplied by a factor and split
equally, so every contributed token returns less than one to its giver but more than one
to the group. The individual incentive is therefore to contribute nothing and keep the
endowment while sharing in what others provide; the collective optimum is for everyone to
contribute everything. The distance between those two is the game's measure of
cooperation.

Human players do not behave as the selfish prediction says. The experimental record,
surveyed by Ledyard~\cite{ledyard1995}, has contributions starting well above zero
and then decaying over repeated rounds as cooperators tire of being free-ridden; Fehr and
G\"achter~\cite{fehrgachter2000} showed the decay reversing when players can pay to punish
defectors, and cooperation collapsing again once they cannot. The game captures both the
pull toward free-riding and the fragility of the cooperation that resists it, which is why
it is the central probe of this thesis.

\subsection{The Trust Game}
\label{subsec:bg-trust}

The Trust Game~\cite{berg1995trust} reduces cooperation to a two-player exchange. An
investor is given an endowment and may send any part of it to a trustee; the sent amount
is tripled in transit, and the trustee then chooses how much to return. Backward induction
predicts that a self-interested trustee returns nothing and a foreseeing investor
therefore sends nothing, yet Berg, Dickhaut, and McCabe found investors sending
substantial amounts and trustees returning enough to reward them. The two roles separate
what the Public Goods Game blends: the investor's move measures trust, the trustee's
measures whether that trust is honored. Played between agents of different capability, the
game asks a pointed question, whether a stronger trustee repays a weaker investor or keeps
the tripled stake.

\subsection{Why iterated mixed-motive games}
\label{subsec:bg-iterated}

Both games are mixed-motive: cooperation and self-interest pull in different directions,
and neither pure cooperation nor pure defection is forced. That is what makes them probes
rather than puzzles. A game with a dominant cooperative action would say nothing about who
exploits whom, and one with a dominant selfish action would say only that the agents found
it. The behavior worth measuring lives in the tension between the two, where an agent's
disposition and its read of the others decide what it does.

Iteration sharpens the reading. A single round is one decision against noise; ten rounds
of the same game expose a consistent strategy, let reputation and retaliation come into
play, and give exploitation time to accumulate into a payoff gap worth measuring. For a
study of capability asymmetry this counts twice. The behavioral signatures that separate a
capable agent from a weak one, steadiness, the timing of defection, the response to a
partner's moves, are visible only across rounds; and the inequality the thesis sets out to
measure is a quantity that builds over a game rather than within a single move. The games
are old. Putting language-model agents inside them is recent, and is the subject of the
next section.

\section{LLM agents in multi-agent systems}
\label{sec:bg-llm-agents}

A language model plays one of these games through its prompt. It is given the rules,
asked for a move in a fixed format, and its answer is played; repeat across rounds and
across agents, and an economic game becomes a multi-agent simulation in which every player
is a language model. A small set of frameworks has grown up to run such studies. FAIRGAME~\cite{fairgame2025} drives matrix games from a JSON
configuration and templated prompts across several providers; Concordia~\cite{concordia2023}
builds richer generative agents for open-ended social scenarios;
NetworkGames~\cite{networkgames2026} places agents with assigned personalities on
small-world and scale-free graphs.

What this young literature has mostly established is how easily the setup can mislead. An
agent's behavior turns out to depend on things that have nothing to do with the game.
Huynh and colleagues~\cite{huynh2025}, extending FAIRGAME to a three-player Public Goods
Game, found an implicit hierarchy effect: the agent named first in the prompt acted as
though it held higher status, so the order of names moved the outcome. A study of
self-recognition in the iterated Public Goods Game~\cite{aimirror2025} found that an agent
led to believe it was facing copies of itself played differently, which means opponents
must be kept anonymous. The wording of the rules matters too; Lore and
Heydari~\cite{loreheydari2023} found language-model contributions in the Public Goods Game
moving with the framing of the prompt. These are not incidental bugs but the working
ground rules of the field, and the design of this thesis follows them: anonymous,
short-lived agent identities, and other-player histories shuffled so that position carries
no meaning (Chapter~\ref{ch:methodology}).

Beneath the cautions, the question this thesis takes up is already stirring. When models
of different sizes are placed in the same game, the smaller ones behave differently: one
study of trust behavior found smaller models tracking human play less closely than larger
ones~\cite{canllmtrust2024}. That a model's capability might shape how it fares against
others, and not only how human-like it looks, is the thread the next section follows into
the work closest to this one.

\section{Capability heterogeneity and cooperation}
\label{sec:bg-capability}

The frameworks of the previous section mostly study agents that are copies of one another,
or that differ in personality. This thesis varies a different thing: capability. What
happens when a strong model and a weak one are put in the same game is the question its
first finding answers, and a small body of work has circled it.

The nearest is a 2026 study titled \emph{More Capable, Less
Cooperative}~\cite{morecapable2026}. In a multi-agent setting where helping is free
and the agents are told to cooperate, it found that capability does not predict cooperation: a
stronger model drove worse collective performance than a weaker one, and several capable models
withheld information they would have lost nothing by sharing. The direction of that result,
capability coming loose from cooperative advantage, is the one this thesis also reports. The two
are best read as complementary, and they differ in instructive ways. That work studies costless
coordination, where the agents share an objective and helping carries no price; this one studies
mixed-motive games, where cooperation is costly and can be exploited. That work runs models in
separate groups; this one mixes the two tiers in a single game and varies their \emph{composition},
from a lone strong agent among the weak to the reverse. It does not measure welfare inequality or
put a number on deception; this thesis does both, and adds a lever the costless setting has no room
for, a communication condition whose effect on the capability gap is the headline of
Chapter~\ref{ch:results}. The two even point in opposite behavioral directions, capable agents
cooperating too little where it is free and too much where it is costly, which is the more telling
for it: capability comes loose from cooperative advantage on both sides.

The observation is older than the framing. Two years earlier, in a study of multi-party
negotiation among language-model stakeholders, Abdelnabi and colleagues~\cite{abdelnabi2023}
noted in passing that
GPT-4 agents earned more than GPT-3.5 agents given the same role in a mixed population, and
that this hinted at questions of fairness and disparity. They did not follow the thread.
That single sentence is, in effect, the hypothesis of this thesis written before it; the
work here turns the aside into a measurement. Related asymmetries surface elsewhere: only
the stronger models in a negotiation self-play study improved from
feedback~\cite{fuetal2023}, and the trust-behavior result of the previous section, smaller
models tracking human play less closely~\cite{canllmtrust2024}, leans the same way. What
none of these did is hold capability as the one thing that varies across a controlled
population and watch what it does to who earns and who is exploited. That is the gap this
thesis fills.

\section{Measuring deception in LLM dialogue}
\label{sec:bg-deception}

If capability is the first axis this thesis measures, deception is the second, and the
harder to pin down. Deception is a gap between what an agent says and what it does, and the
difficulty is turning that gap into a number that means the same thing across agents and
games. Recent work offers several routes, and the choice among them is itself a
methodological question this thesis takes up.

One route compares an agent's \emph{stated} preferences against its \emph{revealed} ones. A
2025 study asked whether language models are consistent between the two and measured the
divergence with a Kullback--Leibler term in single-turn, forced-choice
prompts~\cite{alignmentrevisited2025}. That is the seed of the measure used here, lifted
from one isolated choice to a multi-turn game in which an agent states a policy in a message
and then reveals one in its move. A second route runs deception through a listener: a
multi-turn study scored it by the shift a message produces in another agent's
beliefs~\cite{deceptivedialogue2025}. A third hands the judgment to a model, using a
detector LLM as a proxy for whether a message was deceptive, as in a study built on the
social-deduction game Mafia~\cite{hiddenplaintext2026}; a fourth presents a formalization of
lying as a per-agent honesty cost~\cite{traitors2025}.

The belief-based route is worth pausing on, because it marks where this thesis parts from
the others. Measuring deception by its effect on a listener requires that the message reach
the listener and move its beliefs. The measure used here does not: it compares an agent's
own stated policy, read from its message, against its own action in the same round, asking
only whether one agent's words matched its deeds. That keeps it independent of whether any
message reached another player, a robustness the belief-based measures cannot claim. The
contribution of the thesis is to take the stated-versus-revealed idea seriously enough to
ask which form of it actually tracks deception, a contemporaneous, action-by-action
surprisal set against a global policy divergence, and to report that comparison as a result
rather than assume its answer (Chapters~\ref{ch:methodology} and~\ref{ch:results}).

\section{Positioning: the gap this work fills}
\label{sec:bg-positioning}

The two threads of this chapter, capability and deception, meet at a gap the existing work
leaves open. Each piece of it has been touched; none of it joined. Capability asymmetry has
been seen in coordination tasks~\cite{morecapable2026} and noted in passing in
negotiation~\cite{abdelnabi2023}; heterogeneity has been studied as
personality~\cite{networkgames2026}; deception has been measured through a listener's beliefs
and through detector models~\cite{deceptivedialogue2025, hiddenplaintext2026}. What no single
study does is hold capability as the variable, vary the make-up of the population around it,
and read both the inequality and the deception that follow, across more than one game. That
join is what this thesis supplies, in four parts worth stating plainly.

The first is population composition. Prior comparisons are pairwise, one model against
another; this study fixes the two capability tiers and varies their proportions across a
sixteen-cell matrix, from a lone strong agent in a weak crowd to a lone weak agent among the
strong. Exploitation is a property of a population rather than a pair, and only a population
can show whether a lone capable agent dominates or is ganged up on.

The second is communication, and it has to be stated precisely, because the apparatus is not
what the word suggests. The communication condition asks each agent to make a public pledge
before it acts; it does not pass that pledge to anyone. What the study isolates is therefore
not an exchange of messages but the effect of being asked to commit in public, and its finding
is that this framing alone moves the capability gap (Chapter~\ref{ch:results}). That a public
commitment shifts behavior with no audience to receive it is a narrower and more surprising
claim than ``communication helps,'' and it is the one the design can actually support.

The third is cross-game scope. A result from a single game can be a fact about that game;
running the same populations through a group dilemma and a two-player exchange asks whether an
effect generalizes or belongs to its setting. The fourth is the deception measure of the
previous section, a contemporaneous surprisal validated against human labels and set against a
global policy divergence that the thesis keeps as a negative control rather than a finding.
Taken together, the four turn a scatter of separate observations into one controlled question:
when capable and weak agents share an economy, who exploits whom, does being asked to speak
change it, and can the difference be measured. The chapters that follow build the apparatus to
ask it and report what it answered.
